\documentclass[aspectratio=169,t,13pt,usenames,dvipsnames]{beamer}
\usetheme[lily]{PaloAlto}
\usecolortheme{dolphin}
\usepackage{mathtools}
\usepackage{lmodern}
\usepackage{tikz}
\usepackage{pgfplots}
\usepackage{tikz-3dplot}
\usepackage{adjustbox}
\usetikzlibrary{quantikz}
\usepackage{xcolor}
\usepackage{bm}
\pgfplotsset{compat=newest}
\logo{%
\stackbox[c][c]{\includegraphics[height=1.15cm]{logo.pdf}\\[-3pt]{\textcopyright~2021}\\{Ron K. Cytron}}}

% \logo{\includegraphics[height=2.5cm]{logo.pdf}}
\def\ColorR#1{\textcolor{BrickRed}{#1}}
\def\ColorG#1{\textcolor{\RCthree}{#1}}
\def\ColorB#1{\textcolor{\RCtwo}{#1}}
\definecolor{links}{HTML}{2E8B57}
\hypersetup{colorlinks,linkcolor=,urlcolor=links}
\def\LinkArrow#1{%
\mbox{$\,$\href{#1}{\mbox{\vrule width 0pt height 0ex depth 0pt$\mapsto$}}}$\,$}
\def\WLOG{WLOG}
\def\True{\mbox{\texttt{True}}}
\def\False{\mbox{\texttt{False}}}
\def\Zero{\mbox{$0$}}
\def\One{\mbox{$1$}}
\def\Not#1{%
\ensuremath{\Overline{\ensuremath{#1}}}}
\def\Xor#1#2{\ensuremath{#1\oplus #2}}
\def\Nand#1#2{%
\ensuremath{\mbox{Nand}(#1,#2)}}
\def\And#1#2{%
\ensuremath{#1 \wedge #2}}
\def\Or#1#2{%
\ensuremath{#1 \vee #2}}
\def\Conj#1{%
\ensuremath{#1^{\star}}}
\def\Star#1{%
\Conj{#1}}
\def\EmptyString{\ensuremath{\lambda}}
\def\Mag#1{%
\ensuremath{{|\,#1\,|}}}
\def\Prob#1{%
\ensuremath{\Mag{#1}^{2}}}
\def\CNOT#1#2{%
\ensuremath{\NamedGate{CNOT}(#1,#2)}}
\def\CNOTMatrix{%
\begin{pmatrix*}[r]
1 & 0 & 0 & 0 \\
0 & 1 & 0 & 0 \\
0 & 0 & 0 & 1 \\
0 & 0 & 1 & 0\end{pmatrix*}%
}
\def\CCNOT#1#2#3{%
\ensuremath{\mbox{\NamedGate{CNOT}}(#1,#2,#3)}}
\long\def\ToDo#1{\smash{\textcolor{YellowOrange}{#1}}}
\def\ApplyGate#1#2{%
\ensuremath{#1\,#2}}
\def\Inverse#1{\ensuremath{{#1}^{-1}}}
\def\Quote#1{``#1''}
\def\NamedGate#1{%
\mbox{\bfseries{\textrm{#1}}}}
\def\Identity{\NamedGate{I}}
\def\Hadamard{\NamedGate{H}}
\def\PauliX{\NamedGate{X}}
\def\PauliY{\NamedGate{Y}}
\def\PauliZ{\NamedGate{Z}}
\def\SQBG#1#2#3#4#5{%
\ensuremath{%
#1 \begin{pmatrix*}[r] #2 & #3 \\ #4 & #5\end{pmatrix*}}}
\def\HMatrix{%
\SQBG{\RootTwo{}}{1}{1}{1}{-1}}
\def\XMatrix{%
\SQBG{\relax}{0}{1}{1}{0}}
\def\ZMatrix{%
\SQBG{\relax}{1}{0}{0}{-1}}
\def\YMatrix{%
\SQBG{\relax}{0}{-i}{i}{0}}
\def\IMatrix{%
\SQBG{\relax}{1}{0}{0}{1}}
\def\Co#1{\mathbb{C}}
\def\Implies#1#2{\ensuremath{#1 \rightarrow\ #2}}
\def\Domain#1{\ensuremath{\mathcal{#1}}}
\def\CompClass#1{\Domain{#1}}
\def\SpecialX#1{\ensuremath{#1^{\star}}}
\def\Forall#1#2{\ensuremath{\forall\ #1, #2}}
\def\Degrees#1{\ensuremath{#1^{\circ}}}
\def\RootTwo{\ensuremath{\frac{1}{\sqrt{2}}}}
\def\TwoSupOp#1#2#3{%
\ensuremath{\RootTwo{}\left(#1 #3 #2\right)}}
\def\TwoSup#1#2{%
   \TwoSupOp{\ket{#1}}{\ket{#2}}{+}}
\def\QPlus{\TwoSupOp{\QZero}{\QOne}{+}}
\def\QMinus{\TwoSupOp{\QZero}{\QOne}{-}}
\def\Set#1{\ensuremath{\left\{\,#1\,\right\}}}
\def\BigSkip{%

\bigskip

}

\def\MedSkip{%

\medskip

}
\def\SmallSkip{%

\smallskip

}
\def\SineWave#1{%
{%
\pgfmathsetmacro{\Xone}{#1}
\pgfmathsetmacro{\Xtwo}{-(#1)}
\draw (0,0) sin(1,\Xone) cos (2,0) sin(3,\Xtwo) cos (4,0);
}
}
\newcommand{\tolstrut}{%
  \vrule height\dimexpr\fontcharht\font`\A+.1ex\relax width 0pt\relax
}

\DeclareRobustCommand{\Overline}[1]{%
  \ensuremath{\overline{\raisebox{0pt}[1.07\height]{#1}}}%
}
\newenvironment{TIKZP}[1][scale=1.0]{%
\adjustbox{valign=t}\bgroup
\begin{tikzpicture}[#1]
}{%
\end{tikzpicture}
\egroup
}
\long\def\Definition#1{%
\begin{block}{Definition}
#1
\end{block}}
\long\def\Postulate#1#2{%
\begin{block}{Postulate}
\begin{description} \item[#1] #2\end{description}
\end{block}}
\long\def\TwoUnequalColumns#1#2#3#4{%
\begin{columns}%
\begin{column}{#1}#3\end{column}%
\begin{column}{#2}#4\end{column}%
\end{columns}%
}
\long\def\ThreeUnequalColumns#1#2#3#4#5#6{%
\begin{columns}%
\begin{column}{#1}#4\end{column}%
\begin{column}{#2}#5\end{column}%
\begin{column}{#3}#6\end{column}%
\end{columns}%
}
\long\def\ThreeColumns#1#2#3{%
\ThreeUnequalColumns{0.33\textwidth}{0.33\textwidth}{0.33\textwidth}{#1}{#2}{#3}%
}
\long\def\TwoColumns#1#2{%
\TwoUnequalColumns{0.5\textwidth}{0.5\textwidth}{#1}{#2}%
}
\long\def\OnlyRemark#1#2{%
\only<#1>{\Remark{#2}}}

\long\def\Remark#1{%
\begin{block}{Remark}
#1
\end{block}}
\long\def\Example#1{%
\begin{example} #1\end{example}}
\def\VV{\textit{vice versa}}
\def\QZero{\ket{\hbox to 0.7em{\hss \ensuremath{0}\hss}}}
\def\QOne{\ket{\hbox to 0.7em{\hss \ensuremath{1}\hss}}}
\def\QState#1{\ket{\QName{#1}}}
\def\QName#1{\ensuremath{\psi_{#1}}}
\def\TensOp{\ensuremath{\bm{\otimes}}}
\def\TensProd#1#2{\ensuremath{#1 \TensOp{} #2}}
\def\TensSupProd#1#2{\ensuremath{#1^{\bm{\otimes} #2}}}
\def\PZero{\SQB{1}{0}}
\def\POne{\SQB{0}{1}}
\def\PPlus{\ensuremath{\RootTwo{}\SQB{1}{1}}}
\def\PMinus{\ensuremath{\RootTwo{}\SQB{1}{-1}}}
\def\PHad{\ensuremath\frac{1}{\sqrt{2}}\SQB{1}{1}}
\def\PHadm{\ensuremath\frac{1}{\sqrt{2}}\SQB{1}{-1}}
\def\Intersect#1#2{\ensuremath{#1 \cap #2}}
\def\Union#1#2{\ensuremath{#1 \cup #2}}
\def\Set#1{\ensuremath{\left\{\,#1\,\right\}}}
\def\BinAlph#1{\ensuremath{\Set{0,1}^{#1}}}
\def\AllBits#1#2{\ensuremath{#1\in \BinAlph{#2}}}
\def\DQB#1#2#3#4{%
\ensuremath{\begin{pmatrix*}[r] #1 \\ #2 \\ #3 \\ #4\end{pmatrix*}}}
\def\QQB#1#2#3#4#5#6#7#8{%
\ensuremath{\begin{pmatrix*}[r] #1 \\ #2 \\ #3 \\ #4 \\ #5 \\ #6 \\ #7 \\ #8\end{pmatrix*}}}
\def\SQB#1#2{%
\ensuremath{\begin{pmatrix*}[r] #1 \\ #2\end{pmatrix*}}}
\def\CQB#1#2{%
\ensuremath{\begin{pmatrix*}[r] #1 #2\end{pmatrix*}}}
\def\CDQB#1#2#3#4{%
\ensuremath{\begin{pmatrix*}[r] #1 #2 #3 #4\end{pmatrix*}}}
\def\TZPointDiam{2pt}
\def\TZPoint#1#2#3{%
\draw[fill=black] (#1) circle (\TZPointDiam) node[#3] {#2} ;}
\def\UnitComplexCircle{%
\draw [<->] (-1.5, 0) -- (1.5, 0)  node[right] {$\Re{}$} ;
   \draw [<->] (0,-1.5) -- (0, 1.5) node[above] {$\Im{}$} ;

   \draw (0, 0) circle (1) ;
   }
   \def\TZPEast{\TZPoint{1,0}{$+1$}{above right}}
   \def\TZPNorth{\TZPoint{0,1}{$\NiceI{}$}{above right}}
   \def\TZPWest{\TZPoint{-1,0}{$-1$}{above left}}
   \def\TZPSouth{\TZPoint{0,-1}{$-\NiceI{}$}{below right}}
\def\TZText#1#2#3{%
  \draw (#1) node [#3] {#2};
}
\def\Qiskit{\texttt{qiskit}}
\def\TwoToThe#1{\ensuremath{2^{#1}}}
\def\Exp#1{\ensuremath{e^{#1}}}
\def\NiceI{\ensuremath{\mathfrak{i}}}
\def\ExpPhase#1{\Exp{\,\NiceI\,#1}}
\def\ExpNegPhase#1{\Exp{-\NiceI\,#1}}
%%
%% Not sure this is used
%%
\def\PhaseLabel#1{\phase{\fbox{#1}}}
%%
\def\ColHead#1{\begin{center}\bf #1\end{center}}
\def\Complex#1#2{\ComplexGen{#1}{#2}{+}}
\def\ComplexDiff#1#2{\ComplexGen{#1}{#2}{-}}
\def\ImaginaryM#1{\ComplexDiff{\relax}{#1}}
\def\ComplexGen#1#2#3{%
\ensuremath{#1 #3 i\relax #2}}
\def\Polar#1#2{\ensuremath{#1\hbox to 0.1pt{\hss}\Exp{i\relax#2}}}
%%
%% #1 -- lower left x
%% #2 -- lower left y
%% #3 -- upper right x
%% #4 -- upper right y
%% #5 -- num lines
\def\PFilter#1#2#3#4#5{

    \draw (#1,#2) -- (#3,#2) -- (#3,#4) -- (#1,#4) -- cycle;
    \foreach \i in {1,...,#5}
    {
        \pgfmathsetmacro{\PFilterInc}{#1 + (#3-#1) / #5 * \i}
        % \edef\PFilterInc{\pgfmathresult}
        \draw (\PFilterInc, #2) -- (\PFilterInc, #4);
    }
}
%% #1 - x
%% #2 - y
%% #3 - degrees
%% #4 - options
%% #5 - radius
\def\RadiantArrowsR#1#2#3{%
\foreach \p in {0, #1,...,360} {
   \draw[#2] (0,0)  -- (\p:#3) ;
}
}
\def\RadiantArrows#1#2{%
   \RadiantArrowsR{#1}{#2}{1}
}
\newcommand\RotateTheta[1][scale=1.0]{%
\SquareOutline{}
\begin{scope}[#1]
\fill[white] (0.5,0.5) circle (0.5);
\draw (0.5,0.5) circle (0.5);
\draw[->] (1.0,0.5) arc(0:45:0.5);
\draw (0.5,0.5) node {$\theta$};
\end{scope}
}
\def\SquareOutline{%
\path (0,0) rectangle (1,1);%
}%
\newcommand\Qif[1][scale=1.0]{%
\begin{scope}[#1]
\SquareOutline{}
\clip (0,0) rectangle(1,1);
\RotateAroundCenter{45}{%
\begin{scope}
\draw (0.15,0.15) rectangle(0.85,0.85);
\end{scope}}
\end{scope}
}
\newcommand\LightSource[1][scale=1.0]{%
\begin{scope}[#1]
\SquareOutline{}
\PhotonGun{}
\begin{scope}[shift={(0.75,0.5)}]
    \RadiantArrowsR{22.5}{color=red}{0.25}
\end{scope}
\end{scope}
}
\def\PhotonGun{%
\fill (0,0.375) rectangle (0.5, 0.625);
}
\def\HPolar{%
\begin{scope}[color=red]
\SquareOutline{}
\draw[thick] (0.5,0.4) -- ++(0.35,0)
             ++(-0.35,0.1) -- ++(0.35,0)
             ++(-0.35,0.1) -- ++(0.35,0);
% \fill (0.90,0.5) circle (0.05);
\end{scope}
}
\def\VPolar{%
\begin{scope}[color=red]
\draw[thick] (0.6,0.375) -- ++(0,.250)
             ++(0.1,-.250) -- ++(0,.250)
             ++(0.1,-.250) -- ++(0,.250);
% \fill (0.95,0.5) circle (0.05);
\end{scope}
}
\newcommand\HPolarizedLightSource[1][scale=1.0]{%
\begin{scope}[#1]
\SquareOutline{}
\PhotonGun{}
\HPolar{}
\end{scope}
}
\newcommand\VPolarizedLightSource[1][scale=1.0]{%
\begin{scope}[#1]
\SquareOutline{}
\PhotonGun{}
\VPolar{}
\end{scope}
}
\newcommand\EVBomb[1][scale=1.0]{%
\begin{scope}[#1]
\fill (0.5,0.5) circle(0.5);
\end{scope}%
}
\newcommand\EVBoom[1][scale=1.0]{%
\begin{scope}[#1]
\fill[white] (0.5,0.5) circle(0.5);
\draw[red,thick] (0.5,0.5) circle(0.5);
\clip (0,0) rectangle(1,1);
\begin{scope}
\node[scale=0.35,xshift=28pt,yshift=28pt,starburst,draw,red,fill=orange,line width=1.5pt]
{BOOM!};
\end{scope}
\end{scope}
}
\newcommand\Measurement[1][scale=1.0]{%
\begin{scope}[#1]
\draw[color=white] (0,0) rectangle(1,1);
\begin{scope}
\clip (0.25,0.25) rectangle (1,1);
\draw (0,0) circle(0.9);
\draw (0,0) circle(1.0);
\end{scope}
\draw[thick] (0,0) --(0,1);
\draw[->] (0.25,0.25) -- (0.9,0.9);
\end{scope}
}
\def\TRectangle#1#2#3#4#5#6{%
\begin{scope}[#6]
\draw [fill=#5] (#1,#2) rectangle (#3,#4);
\end{scope}
}
%% #1 -- llx
%% #2 -- lly
%% #3 -- urx
%% #4 -- ury
%% #5 -- angle
\newcommand\Mirror[1][rotate=0]{%
\clip (0,0) rectangle(1,1);
\begin{scope}[#1]
   \SquareOutline{}
   \TRectangle{0}{.4375}{1}{.5625}{Gray}{rotate=0}
\end{scope}
}
\newcommand\BeamSplitter[1][scale=1.0]{%
\begin{scope}[#1]
    \SquareOutline{}
    \TRectangle{0}{.375}{1}{.625}{SkyBlue}{rotate=0}
\end{scope}
}
\def\MZ{%
\LightSource{}
\Shift{2}{0}{\RotateAroundCenter{-45}{\BeamSplitter{}}}
\Shift{2}{-2}{\RotateAroundCenter{-45}{\Mirror{}}}
\Shift{6}{0}{\RotateAroundCenter{-45}{\Mirror{}}}
\Shift{6}{-2}{\RotateAroundCenter{-45}{\BeamSplitter{}}}
\Shift{8}{-2}{\Measurement[color=orange]{}}
\Shift{6}{-4}{\RotateAroundCenter{-90}{\Measurement[color=\RCtwo]{}}}
}
\def\Shift#1#2#3{%
\begin{scope}[shift={(#1,#2)}] #3\end{scope}
}
\def\RotateAroundCenter#1#2{%
\begin{scope}[rotate around={#1:(0.5,0.5)}]
  #2
\end{scope}
}
\def\Vskip#1{\mbox{}\vskip #1\mbox{}}
\def\Hskip#1{\mbox{}\hskip #1\mbox{}}
\author{Ron K.~Cytron}
\institute{Washington University\\Saint Louis, Missouri\\[2em] \textcopyright~2021 All rights reserved by the author}
\setbeamertemplate{footline}[frame number]{}
%\setbeamertemplate{footline}[text line]{%
%  \parbox{\linewidth}{\vspace*{-8pt}\textcopyright~2021\hfill\insertshortauthor\hfill\insertpagenumber}}
%% #1 -- lecture number
%% #2 -- lecture title
%% #3 -- lecture subtitle
%% #4 -- lecture label
\def\SetTitle#1#2#3#4{%
   \title{#2}%
   \subtitle{#3}%
   \date{\today}%
   %\begin{frame}\maketitle\end{frame}%
   \frame{\titlepage}%
   \lecture[#1]{#2: #3}{lec:#4}%
}
\def\Kaye{\href{https://dl.acm.org/doi/10.5555/1206629}{Kaye}}
\def\MikeIke{\href{https://dl.acm.org/doi/10.5555/1972505}{Nielson~and~Chuang}}

\newcounter{ProtocolDialogStep}
\newenvironment{ProtocolDialog}[3]{%
\def\Incr{\stepcounter{ProtocolDialogStep}}
\def\Ref##1{\Incr{}\ThreeUnequalColumns{#1}{#2}{#3}{\arabic{ProtocolDialogStep}. ##1}{\relax}{\relax}}%
\def\Alice##1{\Incr{}\ThreeUnequalColumns{#1}{#2}{#3}{\relax}{\arabic{ProtocolDialogStep}. ##1}{\relax}}%
\def\All##1##2##3{%
\Incr{}
\ThreeUnequalColumns{#1}{#2}{#3}{\arabic{ProtocolDialogStep}. ##1}{##2}{##3}
}
\def\Bob##1{\Incr{}\ThreeUnequalColumns{#1}{#2}{#3}{\relax}{\relax}{
\arabic{ProtocolDialogStep}. ##1}}%
\setcounter{ProtocolDialogStep}{0}%
\ThreeUnequalColumns{#1}{#2}{#3}{\textbf{Ref}}{\textbf{Alice}}{\textbf{Bob}}
}{%
}
%% #1 -- options
%% #2 -- width
%% #3 -- height
%% #4 -- signals
\newenvironment{GateBox}[4][scale=1.0]{%
\edef\Signals{#4}%
\edef\Height{#3}%
\edef\Width{#2}%
\pgfmathsetmacro{\Max}{\Signals}%
\pgfmathsetmacro{\Vsep}{\Height / \Max}%
\pgfmathsetmacro{\Wlen}{\Width / 5}%
\def\ComputeMyY##1{%
   \pgfmathsetmacro{\MyY}{\Height - 0.5 * \Vsep - ##1*\Vsep}}%
\def\Input##1##2{%
   \pgfmathsetmacro{\Xone}{-\Wlen}
   \pgfmathsetmacro{\Xtwo}{0}
   \pgfmathsetmacro{\MyY}{\Height - 0.5 * \Vsep - ##1*\Vsep}
   \draw[->] (\Xone,\MyY) node[left] {##2} -- (\Xtwo,\MyY);
}%
\def\Output##1##2{%
   \pgfmathsetmacro{\Xone}{\Width}
   \pgfmathsetmacro{\Xtwo}{\Width + \Wlen}
   \pgfmathsetmacro{\MyY}{\Height - 0.5 * \Vsep - ##1*\Vsep}
   \draw[->] (\Xone,\MyY) -- (\Xtwo,\MyY) node[right] {##2};
}%
\def\BoxLabel##1{%
\node[draw, fit={(0,0) (\Width,\Height)}, inner sep=0pt, label=center:\mbox{##1}] (A) {};
}%
\begin{TIKZP}[#1]
   \draw (0,0) rectangle(#2,#3);
}{%
\end{TIKZP}}
%%
%% 1 options
%% 2 width
%% 3 height
%% 4 number of qubits
%%
\newenvironment{DotPBox}[4][scale=1.0]{%
\edef\Qubits{#4}%
\edef\Height{#3}%
\edef\Width{#2}%
\pgfmathsetmacro{\Max}{\Qubits}%
\pgfmathsetmacro{\Hsep}{\Width / \Max}%
\def\Qbit##1##2{%
   \pgfmathsetmacro{\MyX}{##1*\Hsep}
   \draw (\MyX,0) rectangle ++(\Hsep,\Height) node[pos=0.5] {\ensuremath{##2}};
}%
\begin{TIKZP}[#1]
}{%
\end{TIKZP}}
\def\PhiTheta#1#2{%
\mbox{\ensuremath{\textcolor{\RCone}{\phi = #1}, \textcolor{\RCtwo}{\theta = #2}}}}
%
% #2 -- orcacle name
% #3 -- upper left input
% #4 -- lower left input
% #5 -- upper right output
% #6 -- lower right output
%
\newenvironment{Oracle}[6][scale=1.0]{%
\begin{GateBox}[#1]{2.5}{1.5}{2}
\BoxLabel{#2}
\Input{0}{#3}
\Input{1}{#4}
\Output{0}{#5}
\Output{1}{#6}
}{%
\end{GateBox}
}
\newenvironment{TTable}[3]{%
\def\TRow##1##2##3{%
##1 & ##2 & ##3 \\}
\begin{tabular}{c|cc}
\TRow{#1}{#2}{#3}\hline
}{%
\end{tabular}
}
%%
%%  1 name
%%  2 one x
%%  3 one f(x)
%%  4 another x
%%  5 another f(x)
\def\SBitTable#1#2#3#4#5{%
\begin{center}
    \begin{TTable}{#1}{$x$}{$f(x)$}
    \TRow{ }{#2}{#3}
    \TRow{ }{#4}{#5}
    \end{TTable}
\end{center}}
%%
%%
\newenvironment{UTable}[6]{%
\def\URow##1##2##3##4##5##6{%
##1 & ##2 & ##3  & ##4 & ##5 & ##6\\}
\begin{tabular}{c|cc|c|cc}
& \multicolumn{2}{|c|}{Inputs} &
& \multicolumn{2}{|c}{Outputs} \\
\URow{#1}{#2}{#3}{#4}{#5}{#6}\hline
}{%
\end{tabular}
}
%%
%%  1 name
%%  2 one x
%%  3 one f(x)
%%  4 another x
%%  5 another f(x)
\def\SBitTable#1#2#3#4#5{%
\begin{center}
    \begin{TTable}{#1}{$x$}{$f(x)$}
    \TRow{ }{#2}{#3}
    \TRow{ }{#4}{#5}
    \end{TTable}
\end{center}}

\newenvironment{BlochSphere}[3][scale=1.0]{
    \bgroup%
    \def\BlochSpherePhiTheta{%
    \tdplotsetrotatedcoords{0}{0}{0};O
        \path[tdplot_rotated_coords] (1.0,0,0) arc (0:115:1.0) node(phi) {};
        \draw[tdplot_rotated_coords,thick,\RCone] (0.3,0,0) arc (0:115:0.3);
        \draw[tdplot_rotated_coords,\RCone,thick,dashed] (0,0,0) -- (phi);
        \draw[\RCone] (0.4,0.4,0) node {$\phi$};
        
        \tdplotsetrotatedcoords{25}{-90}{0};
        \path[tdplot_rotated_coords] (1,0,0) arc (0:45:1) node(theta) {};
        \draw[tdplot_rotated_coords,thick,\RCtwo] (0.4,0,0) arc (0:45:0.4);
        \draw[tdplot_rotated_coords,\RCtwo,thick] (0,0,0) -- (theta);
        \draw[\RCtwo] (-0.4,0,0.3) node {$\theta$};
        
        
    }
    \def\TZPointDiam{0.7pt}%
    \def\ShowAxes{%
    \draw[dashed, tdplot_rotated_coords, black] (1,0,0) node {{\Large\bf x}} (0,1,0) node {{\Large\bf y}} (0,0,1) node {{\Large\bf z}};
    }%
    \tdplotsetmaincoords{#2}{#3}%
    \begin{tikzpicture}[tdplot_main_coords, #1]
        
        
        \shade[ball color = lightgray, opacity = 0.5] (0,0,0) circle (1cm);
        
        %%
        %% The equator
        %%
        \tdplotsetrotatedcoords{0}{0}{0};
        \draw[dashed, tdplot_rotated_coords, gray] (0,0,0) circle (1);
        
        %%
        %% Arc in the x-z plane
        %%
        \tdplotsetrotatedcoords{90}{90}{90};
        \draw[dashed, tdplot_rotated_coords, gray] (1,0,0) arc (0:180:1);
        
        %%
        %% Arc in the y-z plane
        %%
        \tdplotsetrotatedcoords{0}{90}{90};
        \draw[dashed, tdplot_rotated_coords, gray] (1,0,0) arc (0:180:1);
        
        %\tdplotsetrotatedcoords{90}{90}{90};
       % \draw[dashed, tdplot_rotated_coords, black] (1,0,0) node {{\Large\bf x}} (0,1,0) node {{\Large\bf y}} (0,0,1) node {{\Large\bf z}}; 
        
        \draw[dashed, gray] (0,0,0) -- (-1,0,0);
        \draw[dashed, gray] (0,0,0) -- (0,-1,0);
        \draw[dashed, gray] (0,0,0) -- (0,0,-1);
        
        \draw[-stealth] (0,0,0) -- (1.30,0,0);
        \draw (1.45,0,0) node {$x$};
        \draw[-stealth] (0,0,0) -- (0,1.30,0);
        \draw(0,1.45,0) node {$y$};
        \draw[-stealth] (0,0,0) -- (0,0,1.30);
        \draw (0,0,1.45) node {$z$};
        
}{%
    \end{tikzpicture}
    \tdplotsetmaincoords{0}{0}
    \egroup
}

\long\def\Exercise#1#2{%
\relax}
\newcommand\QEDsym[1][2em]{\hbox to #1{\hss\ensuremath{\square}}}
\def\QED{\begin{flushright}\ensuremath{\square}\end{flushright}}
\def\UGateA{%
\ensuremath{\cos(\theta/2)}}
\def\UGateB{%
\ensuremath{-\ExpPhase{\lambda}\sin(\theta/2)}}
\def\UGateC{%
\ensuremath{\ExpPhase{\phi}\sin(\theta/2)}}
\def\UGateD{%
\ensuremath{\ExpPhase{(\phi+\lambda)}\cos(\theta/2)}}
\def\UGate{%
\SQBG{\relax}{\UGateA{}}{\UGateB{}}{\UGateC{}}{\UGateD{}}}%
\def\FCirc#1{%
\begin{tikzpicture}[scale=#1]%
\SquareOutline{}%
\fill (0.5,0.5) circle(0.4);%
\end{tikzpicture}%
}%
\def\SCirc{\FCirc{0.3}}%
\newenvironment{BBKey}[0]{%
\def\STD{%
\begin{tikzpicture}[scale=0.35]%
\SquareOutline{}%
\draw[ultra thick,purple] (0,0.5) -- (1,0.5);%
\draw[ultra thick,purple] (0.5,0) -- (0.5,1);%
\end{tikzpicture}}%
\def\HDM{%
\begin{tikzpicture}[scale=0.30]%
\SquareOutline{}%
\draw[ultra thick,purple] (0,0) -- (1,1);%
\draw[ultra thick,purple] (0,1) -- (1,0);%
\end{tikzpicture}}%
\def\BBUp{%
\begin{tikzpicture}[scale=0.30]%
\SquareOutline{}%
\draw[->,thick,purple] (0.5,0) -- (0.5,1);%
\end{tikzpicture}}%
\def\BBRt{%
\begin{tikzpicture}[scale=0.30]%
\SquareOutline{}%
\draw[->,thick,purple] (0,0.5) -- (1,0.5);%
\end{tikzpicture}}%
\def\BBNe{%
\begin{tikzpicture}[scale=0.30]%
\SquareOutline{}%
\draw[->,thick,purple] (0,0) -- (1,1);%
\end{tikzpicture}}%
\def\BBSe{%
\begin{tikzpicture}[scale=0.30]%
\SquareOutline{}%
\draw[->,thick,purple] (0,1) -- (1,0);%
\end{tikzpicture}}%
\def\BBBasis{%
\begin{tabular}{c|c|c}
Basis & 0 & 1 \\
\hline
\STD{} & \BBUp{} & \BBRt{} \\
\HDM{} & \BBNe{} & \BBSe{}
\end{tabular}%
}%
}{%
\relax%
}
\def\CFCirc#1{%
  \textcolor{#1}{\FCirc{0.30}}}
\def\Agree{%
  \textcolor{green}{\FCirc{0.30}}}%
\def\DisAgree{%
   \textcolor{red}{\FCirc{0.30}}}%
\def\MetaAgree{%
   \textcolor{yellow}{\FCirc{0.30}}}%
\newenvironment{EKey}[0]{%
\def\ZZ{%
\ZRot{0}%
}%
\def\ZZPO##1{%
\ZRot{\frac{\pi}{##1}}%
}%
\def\ZZNPO##1{%
\ZRot{\ensuremath{\,-\frac{\pi}{##1}}}%
}%
\def\ZRot##1{%
\ensuremath{Z_{##1}}%
}%
}{%
}
\def\AltAlert#1#2#3{%
\alt<#1>{\alert{#2}}{#3}}
\def\RCone{Rhodamine}
\def\RCtwo{CornflowerBlue}
\def\RCthree{LimeGreen}
\def\RCfour{Tan}
\def\RCfive{Plum}
\def\ColorOne#1{\textcolor{\RCone}{#1}}
\def\ColorTwo#1{\textcolor{\RCtwo}{#1}}
\def\ColorThree#1{\textcolor{\RCthree}{#1}}
\def\ColorFour#1{\textcolor{\RCfour}{#1}}
\def\ColorFive#1{\textcolor{\RCfive}{#1}}
\def\ColorBlack#1{\textcolor{black}{#1}}
\def\NegF#1{\NegOneExp{f(#1)}}
\def\NegOneExp#1{\ensuremath{(-1)^{#1}}}
\def\RootTwoN#1{\ensuremath{\frac{1}{\sqrt{2^{#1}}}}}
%%
%% var, exponent
%%
\def\SumBV#1#2{\OpBV{#1}{#2}{\displaystyle\sum\limits}}
\def\OpBV#1#2#3{\ensuremath{#3_{\AllBits{#1}{#2}}}}
\def\DotP#1#2{\ensuremath{#1\boldsymbol{\cdot}#2}}
\def\SumPH#1#2{\displaystyle\sum\limits_{#1=0}^{2^{#2}-1}}
%
% for hadamard results
%
%% #1 w for H((w>)
%% #2 free variable name
%% #3 n
\newcounter{HTableCounter}
\newenvironment{HTable}[3]{%
%% ##1 v^w
%% ##2 v . w
%% ##3 result
%% ##4 basis vector 000 001 etc.
\def\TE##1##2##3##4{%
\ensuremath{##4} & \ensuremath{##1} & \ensuremath{##2} & \ensuremath{##3 1} \ensuremath{\ket{##4}}}%
\begin{tabular}{c|c|c|r}
\multicolumn{4}{c}{$\sqrt{2^{#3}}$\,\Hadamard(\ket{#1})} \\ 
$v$ & \And{v}{#1} & \DotP{v}{#1} &  \multicolumn{1}{|c}{Term} \\ \hline}{%
\end{tabular}
}

\newcounter{ScrollCounter}
\newcounter{ScrollCounterEnd}
\newcounter{ScrollEOne}
\newcounter{ScrollETwo}
\newcounter{ScrollEThree}
\newcounter{ScrollEFour}
%%
%% start
%%
\newenvironment{ScrollingProof}[1][1]{%
   \setcounter{ScrollCounter}{#1}%
   \setcounter{ScrollEOne}{#1}%
   \setcounter{ScrollETwo}{#1}%
      \stepcounter{ScrollETwo}%
   \setcounter{ScrollEThree}{#1}%
      \stepcounter{ScrollEThree}\stepcounter{ScrollEThree}%
   \setcounter{ScrollEFour}{#1}%
      \stepcounter{ScrollEFour}\stepcounter{ScrollEFour}\stepcounter{ScrollEFour}
   \def\Curr{\value{ScrollCounter}}%
   \def\One{\value{ScrollEOne}}%
   \def\Two{\value{ScrollETwo}}%
   \def\Three{\value{ScrollEThree}}%
   \def\Four{\value{ScrollEFour}}%
   %\def\Memo##1{\gdef\RCScrollMemoText{##1}##1}
   %\def\Recall{\RCScrollMemoText}
   \def\Next##1##2{%
       \visible<\value{ScrollCounter}-##1>{##2}%
       %\def\Last{##2}%
    \stepcounter{ScrollCounter}%
    \stepcounter{ScrollEOne}%
    \stepcounter{ScrollETwo}%
    \stepcounter{ScrollEThree}%
    }%
    \def\Last##1{%
        \gdef\ScrollingLastText{##1}%
        \Next{\Two}{\visible<\value{ScrollETwo}>{\alert{\ensuremath{\big\Uparrow}}\ }##1}}%
}{%
}
\long\def\ScrollProof#1#2#3{%
\only<#1-#2>{%
\begin{ScrollingProof}[#1]
\begin{align*}
#3
\end{align*}
\end{ScrollingProof}}}
\def\RVDots{\raisebox{0.5em}{\ensuremath{\vdots}}}
%%
%% #1 -- the state undergoing hadamard
%% #2 -- the number of bits
%% #3 -- the free variable
\def\NHadamard#1#2#3{%
\SumBV{#3}{#2} \NegOneExp{\DotP{#1}{#3}} \ket{#3}}
\newenvironment{Reasoning}{%
\def\Reason##1##2{%
   \only<##1>{\makebox(0,0)[l]{\alert<##1>{\textit{##2}}}}}%
%%
%% #1 -- the range for overlays
%% #2 -- the color
%% #3 -- x
%% #4 -- y
%% #5 -- width
%% #6 -- height
%%
\def\CallOut##1##2##3##4##5##6{%
    \Reason{##1}{\Highlight{##2}{##3}{##4}{##5}{##6}}}
\vrule width 0pt depth 0.5em height 0.5em}{%
}
\def\Span#1{\ensuremath{{\mbox{span}(#1)}}}
\def\EmptySet{\Set{\,}}
%%
%% Large Binary Point
%%
\def\LBP{\mbox{\Huge .\hskip -1pt}}
%%
%% #1 color
%% #2 x loc
%% #3 y loc
%% #4 width
%% #5 height
\def\Highlight#1#2#3#4#5{%
\begin{tikzpicture}[overlay]
\draw[fill=#1,opacity=0.2] (#2,#3) rectangle ++(#4,#5);
\end{tikzpicture}}
\def\Indent#1{\hbox to #1{\hss}}


% \AtBeginLecture{\frame{\Large \insertlecture}}
%%
%% test
%%
\begin{document}
\SetTitle{30}{Phase estimation}{What separates the superpositions in a state?}{30}

\begin{frame}{Overview}{What will we study here?}

\begin{itemize}[<+->]
    \item Phase has been an elusive aspect of quantum state.
    \item As we have studied, the probability of a measurement outcomes is related to the amplitude of the coefficient present the associated basis state.
    \item We use phase during a quantum computation to establish interference patterns of interest.
    \item While we cannot directly measure phase, we can estimate a constant phase shift between basis states in a superposition.
    \item The study of this material leads directly to quantum Fourier analysis, which in turn is the basis of Shor's famous algorithm.
\end{itemize}
    
\end{frame}

\section*{Phase shift}

\begin{frame}{A return to the complex circle}{The basics we need for phase estimation}
\TwoUnequalColumns{0.6\textwidth}{0.4\textwidth}{%
\adjustbox{valign=t, width=\textwidth}{\begin{TIKZP}
\UnitComplexCircle{}
\visible<1>{%
\TZPEast{}\TZPNorth{}\TZPSouth{}\TZPWest{}}
\visible<2>{%
  \draw[->,thick,\RCtwo] (1,0) arc (0:45:1) coordinate (z) ;
  \draw[thick,dotted,\RCtwo] (0,0) -- (z) ;
  \TZText{0.5,0}{$\theta$}{above}
    \ColorTwo{\TZText{z}{\ExpPhase{\theta}}{above right}}
}
\visible<3>{%
   \TZPoint{1,0}{$\ExpPhase{0}=+1$}{above right}}
\visible<4>{
      \TZPoint{0,1}{$\ExpPhase{\pi/2}=i$}{above right}
      \draw[->,thick,\RCtwo] (1,0) arc (0:90:1) coordinate (z) ;
    \draw[thick,dotted] (0,0) -- (z);}
\visible<5>{%
      \TZPoint{-1,0}{$\ExpPhase{\pi}=-1$}{above left}
      \draw[->,thick,\RCtwo] (1,0) arc (0:180:1) coordinate (z) ;
    \draw[thick,dotted] (0,0) -- (z);
    }
\visible<6>{%
      \TZPoint{0,-1}{$\ExpPhase{3\pi/2}=-i$}{below right}
      \draw[->,thick,\RCtwo] (1,0) arc (0:270:1) coordinate (z) ;
    \draw[thick,dotted] (0,0) -- (z);
    }
\visible<7-8>{%
      \TZPoint{1,0}{$\ExpPhase{2\pi}\visible<8>{=}$}{above right}
       \visible<8>{\TZPoint{1,0}{$\ExpPhase{0}=+1$}{below right}}
      \draw[->,thick,\RCtwo] (1,0) arc (0:359:1) coordinate (z) ;
    \draw[thick,dotted] (0,0) -- (z);
}
\end{TIKZP}}}{%
\begin{itemize}
    \item<1-> Unit points on each axis
    \item<2-> Traveling $\theta$ radians counterclockwise from $+1$
\end{itemize}
}
    
\end{frame}

\begin{frame}{Constant phase shift}{On the complex unit circle}
\TwoUnequalColumns{0.4\textwidth}{0.6\textwidth}{%
\adjustbox{valign=t, width=\textwidth}{\begin{TIKZP}
\UnitComplexCircle{}
\visible<3-8>{%
 \draw[->,thick,\RCone] (1,0) arc (0:40:1) coordinate (z2);
 \draw[thick,dotted] (0,0) -- (z2);
 \visible<4->{\draw[->,thick,\RCone] (1,0) arc (0:60:1) coordinate (z3);
 \draw[thick,dotted] (0,0) -- (z3);}
 \visible<5->{\draw[->,thick,\RCone] (1,0) arc (0:80:1) coordinate (z4);
 \draw[thick,dotted] (0,0) -- (z4);}
 \visible<6->{\draw[->,thick,\RCone] (1,0) arc (0:100:1) coordinate (z5);
 \draw[thick,dotted] (0,0) -- (z5);}
}
\visible<1-8>{%
  \draw[->,thick,\RCtwo] (1,0) arc (0:20:1) coordinate (z) ;
  \draw[thick,dotted,\RCtwo] (0,0) -- (z);
  \visible<1>{\TZText{0.55,0.15}{\mbox{\small $\theta$}}{right}}
  \visible<1>{
    \ColorTwo{\TZText{z}{\ExpPhase{\theta}}{above right}}}
  \visible<2->{
    \ColorTwo{\TZText{z}{\ExpPhase{w\cdot 2\pi}}{above right}}}
}
\visible<9->{%
  \draw[->,ultra thick,\RCthree] (0,0) -- (1,0);
  \visible<9->{\ColorThree{\TZPoint{1,0}{$\ExpPhase{0}=+1$}{below right}}}
  \visible<10->{%
  \draw[->,thick,\RCtwo] (1,0) arc (0:20:1) coordinate (z2);
  \draw[thick,dotted] (0,0) -- (z2);
  \ColorTwo{\TZText{z2}{\ \ExpPhase{w\cdot 2\pi\cdot 1}}{right}}}
  \visible<11->{%
 \draw[->,thick,\RCone] (1,0) arc (0:40:1) coordinate (z3);
 \draw[thick,dotted] (0,0) -- (z3);
 \ColorOne{\TZText{z3}{\ \ExpPhase{w\cdot 2\pi\cdot 2}}{right}}}
 \visible<12->{%
 \draw[->,thick,\RCone] (1,0) arc (0:60:1) coordinate (z4);
 \draw[thick,dotted] (0,0) -- (z4);
 \ColorOne{\TZText{z4}{\ \ExpPhase{w\cdot 2\pi\cdot 3}}{above right}}}
 \visible<13->{%
 \draw[->,thick,\RCone] (1,0) arc (0:80:1) coordinate (z5);
 \draw[thick,dotted] (0,0) -- (z5);
 \ColorOne{\TZText{z5}{\raisebox{4pt}{\ \ExpPhase{w\cdot 2\pi\cdot 4}}}{above right}}}
}
\end{TIKZP}}
}{%
\only<1-7>{%
\begin{itemize}
    \item<1-> The basis states will be related to each other in phase by a constant shift of $\theta$ radians.
    \item<2-> We capture the extent of the rotation by $\ColorTwo{w}, 0\leq \ColorTwo{w} < 1$, which represents a fraction of a complete rotation of $2\pi$ radians.
    \item<3-> Each \ColorOne{subsequent basis vector} has amplitude shifted by $\ColorTwo{w}\cdot 2\pi$ radians.
    \item<7-> We generate as many of these as need to cover the basis vectors, $2^n$ for $n$~qbits.
\end{itemize}}
\only<8->{%
\begin{itemize}
    \item<8-> Global phase consists of a uniform rotation of all points on the complex circle. 
    \item<9-> We can therefore assert that the first point occurs at $+1$.
    \item<10-> Each \ColorOne{subsequent point} is $\ColorTwo{w}\cdot 2\pi$ radians counterclockwise on the circle.
    \item<14-> We can characterize these as
    \[ \ExpPhase{w\cdot 2\pi \cdot y}, 0\leq y < 5\]
\end{itemize}}
}
\end{frame}

\begin{frame}{Periodic states}{definition}

A \emph{periodic state} \QState{} for $n$ qubits can be described as:
\[
\QState{} = \textcolor<4>{\RCtwo}{\RootTwoN{n}}\SumPH{\textcolor<3>{\RCone}{y}}{n} \ExpPhase{2\pi w\textcolor<3>{\RCone}{y}} \ket{\textcolor<3>{\RCone}{y}}
\]
for a constant $w, 0 \leq w \textcolor<2>{\RCthree}{< 1}$.
\begin{itemize}
    \item<2-> Some sources include~$1$ in \textcolor<2>{\RCthree}{the upper bound for~$w$}, but it's the same as~$w=0$.
    \item<3-> The iteration variable \textcolor<3>{\RCone}{$y$} interprets the binary string encoding a basis vector as a number.
    \item<4-> The basis states all have the same magnitude~\textcolor<4>{\RCtwo}{\RootTwoN{}}, so each is equally likely upon measurement.
\end{itemize}
    
\end{frame}

{
\def\V#1#2{\visible<#1->{\alert<#1>{\ExpPhase{2\pi \ColorTwo{\frac{1}{8}}\cdot #2}}}}
\begin{frame}{A column-vector view}{Of a periodic state}
\TwoColumns{%
\adjustbox{valign=t, width=\textwidth}{\begin{TIKZP}
\UnitComplexCircle{}
\visible<2,6->{%
\draw[->,ultra thick,\RCthree] (0,0) -- (1,0);
  \ColorThree{\TZPoint{1,0}{\ExpPhase{0}}{above right}}
}
\visible<3,6->{%
\draw[->,thick,\RCone] (1,0) arc (0:45:1) coordinate (z2);
  \draw[thick,dotted] (0,0) -- (z2);
  \ColorOne{\TZText{z2}{\ \ExpPhase{\frac{\pi}{4}}}{right}}
}
\visible<4,6->{%
\draw[->,thick,\RCone] (1,0) arc (0:90:1) coordinate (z3);
  \draw[thick,dotted] (0,0) -- (z3);
  \ColorOne{\TZText{z3}{\ \ExpPhase{\frac{\pi}{2}}}{above right}}
}
\visible<5,6->{%
\draw[->,thick,\RCone] (1,0) arc (0:135:1) coordinate (z4);
  \draw[thick,dotted] (0,0) -- (z4);
  \ColorOne{\TZText{z4}{\ \ExpPhase{\frac{3\pi}{4}}}{above left}}
}
\end{TIKZP}}
}{%
\ColorTwo{\[w = \frac{1}{8} \]}
\[\RootTwoN{n}\DQB{\V{2}{0}}{\V{3}{1}}{\V{4}{2}}{\V{5}{3}}\]
}
\end{frame}}

\begin{frame}{Some familiar periodic states}{We saw these before but now regard them as taking a trip around the unit complex circle}

\Vskip{-5em}\TwoColumns{%
\begin{align*}
    \ket{+} &= \QPlus{} \\
           \visible<2->{ &=  \RootTwoN{1}\left(\ColorOne{\ExpPhase{0}}\QZero{} \visible<3->{+ \ColorTwo{\ExpPhase{0}}\QOne{}}\right) \\}
        \visible<4->{    &= \RootTwoN{1}\SumPH{y}{1} \ExpPhase{2\pi\cdot \ColorFour{w}\cdot y} \ket{y}, \ColorFour{w=0}}
\end{align*}
}{%
\begin{align*}
    \ket{-} &= \QMinus{} \\
           \visible<5->{ &=  \RootTwoN{1}\left(\ColorOne{\ExpPhase{0}}\QZero{} \visible<6->{+ \ColorTwo{\ExpPhase{\pi}}\QOne{}}\right) \\}
        \visible<7->{    &= \RootTwoN{1}\SumPH{y}{1} \ExpPhase{2\pi\cdot \ColorFour{w}\cdot y} \ket{y}, \ColorFour{w=\frac{1}{2}}}
\end{align*}
}
\Vskip{-4em}\TwoColumns{%
\begin{center}
\adjustbox{valign=t, width=0.6\textwidth}{\begin{TIKZP}
\UnitComplexCircle{}
  \visible<2->{
  \draw[->,ultra thick,\RCone] (0,0) -- (1,0);
  \ColorOne{\TZPoint{1,0}{\ExpPhase{0}}{above right}}}
  \visible<3->{
  \draw[->,ultra thick,\RCtwo] (0,0) -- (1,0);
  \ColorTwo{\TZPoint{1,0}{\ \ \ \ \ \ \ExpPhase{0}}{above right}}}

\end{TIKZP}}
\end{center}
}{%
\begin{center}
\adjustbox{valign=t, width=0.6\textwidth}{\begin{TIKZP}
\UnitComplexCircle{}
  \visible<5->{
  \draw[->,ultra thick,\RCone] (0,0) -- (1,0);
  \ColorOne{\TZPoint{1,0}{\ExpPhase{0}}{above right}}}
  \visible<6->{
  \draw[->,ultra thick,\RCtwo] (0,0) -- (-1,0);
  \ColorTwo{\TZPoint{-1,0}{\ExpPhase{\pi}}{above left}}}

\end{TIKZP}}
\end{center}
}
    
\end{frame}

\begin{frame}{Phase estimation}{Of periodic states}

\begin{itemize}[<+->]
    \item Since $0\leq w < 1$, we can approximate $w$ using the binary fraction
    \[ \LBP{}w_{1}w_{2}\cdots w_{n}\]
    for $n$ bits, where each is~$0$ or~$1$.
    \item For a single qubit to characterize phase, we have
    \begin{description}
        \item[$w=0$]  $\LBP{}w_{1}=.0$ characterizes \ket{+} exactly
        \item[$w=\frac{1}{2}$] $\LBP{}w_{1}=.1$ characterizes \ket{-} exactly
    \end{description}
    \item No extra precision is required to capture \ket{+} or \ket{-}, but other periodic states require more qubits for more precision.
    \item Some periodic states cannot be captured precisely with any number of qubits:
    \begin{itemize}
        \item A state whose period is irrational
        \item A state whose period is rational but whose base-$2$ representation does not terminate
    \end{itemize}
\end{itemize}
    
\end{frame}

\begin{frame}{Phase estimation}{Problem statement}

Given \emph{periodic state} \QState{} for $n$ qubits: 
\[
\QState{} = \RootTwoN{n}\SumPH{y}{n} \ExpPhase{2\pi wy} \ket{y},  0 \leq w < 1
\]
use $n$ qubits to compute an $n$-bit estimation of $w$:
\[
w \approx \LBP{}w_{1}w_{2}\cdots w_{n} 
\]
    
\end{frame}

\begin{frame}{Some intuition}{Using the simplest periodic states}

\begin{itemize}[<+->]
    \item If we know $\QState{}=\ket{+}$ or $\QState{}=\ket{-}$
    \item then we know only a single qubit is needed to represent the phase.
    \item $\Hadamard(\QState{})$ yields the bit we seek:
    \begin{description}
        \item[$\QState{}=\ket{+}$] $\Hadamard(\QState{}) = \ket{0}$
        \item[$\QState{}=\ket{-}$] $\Hadamard(\QState{}) = \ket{1}$
    \end{description}
    \item Measuring the outcome of the \Hadamard{} gate provides $\LBP{}w_{1}$ directly.
    \end{itemize}
    \end{frame}
    
\begin{frame}{And if there are more bits}{Building intuition}
    \TwoColumns{%
If we can isolate each bit after a binary point, then we can recover the other binary digits as well, using the \Hadamard{} gate.  
\SmallSkip{}
This requires:
    \begin{itemize}
        \item<2-> Shifting the binary point right, so that it abuts \ColorTwo{the bit of interest}
        \item<3-> Dropping the \ColorOne{bits to the left}, because we can
        \item<4-> Rotating to eliminate \ColorThree{any non-zero bits} after \ColorTwo{the bit of interest}
        \item<5-> What remains is~$.0$ or~$.1$, and an \Hadamard{} gate can tell us which it is.
    \end{itemize}
}{%
\begin{align*}
\LBP{}\ColorOne{w_{1}w_{2}\cdots w_{k}}\invisible{\LBP{}}\ColorTwo{w_{k+1}}\ColorThree{w_{k+2}w_{k+3}\cdots w_{n}} \\
\visible<2->{\invisible{\LBP{}}\ColorOne{w_{1}w_{2}\cdots w_{k}}\LBP{}\ColorTwo{w_{k+1}}\ColorThree{w_{k+2}w_{k+3}\cdots w_{n}} \\}
\visible<3->{\invisible{\LBP{}\ColorOne{w_{1}w_{2}\cdots w_{k}}}\LBP{}\ColorTwo{w_{k+1}}\ColorThree{w_{k+2}w_{k+3}\cdots w_{n}} \\}
\visible<4->{\invisible{\LBP{}\ColorOne{w_{1}w_{2}\cdots w_{k}}}\LBP{}\ColorTwo{w_{k+1}}\invisible{\ColorThree{w_{k+2}w_{k+3}\cdots w_{n}}} }
%% left off here
\end{align*}
}
    
\end{frame}

\begin{frame}{Example}{With $n=3$ qubits and $w=\frac{5}{8}$}

\Vskip{-3em}\TwoColumns{%
\only<1-7>{%
Finding \ColorThree{$w_3$}
\begin{itemize}
    \item<2-> Shift the binary point right two bits
    \item<4-> Drop the bits to the left
    \item<6-> Applying \Hadamard{}, say we get \ColorThree{1}
\end{itemize}}%
\only<8-15>{%
Finding \ColorTwo{$w_2$}
\begin{itemize}
    \item<8-> Shift the binary point right one bit
    \item<10-> Drop the bit to the left
    \item<12-> Rotate away \ColorThree{$w_3$}=\ColorThree{$1$}
    \item<14-> Applying \Hadamard{}, say we get \ColorTwo{0}
\end{itemize}
}%
\only<16-22>{%
Finding \ColorOne{$w_1$}
\begin{itemize}
    \item<16-> No need to shift the binary point
    \item<17-> Rotate away \ColorThree{$w_3$}=\ColorThree{$1$}
    \item<19-> Must not rotate away \ColorTwo{$w_2$}=\ColorTwo{$0$}
    \item<21-> Applying \Hadamard{}, say we get \ColorOne{1}
\end{itemize}
}%
}{%
\[w=
\visible<1-2,8,16->{\LBP{}}
\visible<1-4,8-10,16->{\ColorOne{w_{1}}}\visible<9-15>{\LBP{}}
\visible<1-4,8-19,23->{\ColorTwo{w_{2}}}\visible<3-7>{\LBP{}}
\visible<1-12,16-17,23->{\ColorThree{w_{3}}}
=
\visible<1-2,8,16->{\LBP{}}
\visible<22->{\ColorOne{1}}\visible<9-15>{\LBP{}}
\visible<15->{\ColorTwo{0}}\visible<3-7>{\LBP{}}
\visible<7->{\ColorThree{1}}
\visible<23->{= \frac{5}{8}}\]
}
    
\end{frame}

\begin{frame}{Some preliminaries}{Results needed for phase estimation}

\begin{itemize}[<+->]
    \item Consider $\ExpPhase{2\pi w}, w=\LBP{}\ColorOne{w_{1}w_{2}\cdots w_{k}}\ColorTwo{w_{k+1}\cdots w_{n}}$ (base 2)
    \item Multiplying $w$ by a power of two shifts the binary point to the right:
    \begin{align*}
       w &= \LBP{}\ColorOne{w_{1}w_{2}\cdots w_{k}}\ColorTwo{\invisible{\LBP{}}w_{k+1}\cdots w_{n}} \\
       \visible<3->{ 2^{k}w &= \ColorOne{\invisible{\LBP{}}w_{1}w_{2}\cdots w_{k}}\LBP{}\ColorTwo{w_{k+1}\cdots w_{n}},\ 0\leq k\leq n}
    \end{align*}
    
    \item<4-> In an exponent, $\ColorOne{w_{1}w_{2}\cdots w_{k}}$ before \LBP{} can be ignored:
    {\Large
    \begin{align*}
        \visible<4->{\ExpPhase{2\pi 2^{k}w} &= \ExpPhase{2\pi \ColorOne{w_{1}w_{2}\cdots w_{k}}\LBP{}\ColorTwo{w_{k+1}\cdots w_{n}}}} \visible<5->{=\ExpPhase{2\pi (\ColorOne{w_{1}w_{2}\cdots w_{k}}+\LBP{}\ColorTwo{w_{k+1}\cdots w_{n}})}\\}
        \visible<6->{&= \ExpPhase{2\pi \ColorOne{w_{1}w_{2}\cdots w_{k}}} \times \ExpPhase{2\pi (\LBP{}\ColorTwo{w_{k+1}\cdots w_{n}})} \\}
       \visible<7->{ &= \ColorOne{1} \times \ExpPhase{2\pi (\LBP{}\ColorTwo{w_{k+1}\cdots w_{n}})}, \mbox{\normalsize because $\ExpPhase{2\pi \ColorOne{N}}=\ColorOne{1}$ for any integer \ColorOne{$N$}}}
    \end{align*}}
\end{itemize}


    
\end{frame}

{
\def\T#1{\TwoSupOp{\QZero}{\ExpPhase{2\pi\left(2^{#1}\W{}\right)}\QOne}{+}}
\def\N{\ColorOne{\ensuremath{n}}}\def\W{\ColorTwo{\ensuremath{w}}}
\begin{frame}{Useful theorem}{Needed for phase estimation, prove later}
\Vskip{-3em}\begin{align*}
 \RootTwoN{\N}\SumPH{y}{\N} \ExpPhase{2\pi \W{}y} \ket{y} =& \T{\N{}-1}  \\
 \TensOp{} & \T{\N{}-2} \\
 \TensOp{} & \cdots \\
 \TensOp{} & \T{\N{}-\N{}} 
\end{align*}
\begin{itemize}
    \item This articulates the state of each of the \N{} qubits separately.
    \item Each successive qubit reveals another binary digit of \W{}.
\end{itemize}
\end{frame}}



{
\def\T#1{\TwoSupOp{\QZero}{\ExpPhase{2\pi\left(2^{#1}\W{}\right)}\QOne}{+}}
\def\N{\ColorOne{\ensuremath{n}}}\def\W{\ColorTwo{\ensuremath{w}}}
\def\S#1#2{\ExpPhase{2\pi\W{}#1}\ket{#2}}
\begin{frame}{How are these expressions different?}{They are two views of the same state}

\TwoColumns{%
\Vskip{-3em}\begin{center}
\adjustbox{scale=0.8,valign=t}{%
\begin{quantikz}
\lstick{\alert<6>{$q_2$}} & \qw & \qw & \rstick[wires=3]{\QState{}=\ket{\alert<2,6>{q_{2}}\alert<2,7>{q_{1}}\alert<2,8>{q_{0}}}} \\
\lstick{\alert<7>{$q_1$}} & \qw & \qw  & \\
\lstick{\alert<8>{$q_0$}} & \qw & \qw  & \\
\end{quantikz}}%
%%
%% force distance at the bottom to keep the second two-columns still
%%
\hrule width 0pt depth 2em
%%
\end{center}}{%
\Vskip{-3em}\only<1-3>{%
\begin{itemize}
    \item<1-> This expression has $2^{n}$ terms.
    \item<2-> Each term represents a particular 3-qubit outcome
    \item<3-> The phases on each successive term are related by \W{}.
\end{itemize}
}%
\only<4->{%
\begin{itemize}
    \item<4-> This expression has $n$ terms.
    \item<5-> Each represents the tensor-factored state of an individual qubit.
    \item<9-> Each exposes more of the binary fraction $w$.
\end{itemize}
}%
}
\Vskip{-5em}\TwoColumns{%
\begin{align*}
\QState{} =& \RootTwoN{3}\SumPH{y}{3} \ExpPhase{2\pi \W{}y} \ket{y} \\[.4em]
=& \S{0}{000} + \S{1}{001} \\
+& \S{2}{010} + \S{3}{011} \\
+& \S{4}{100} + \S{5}{101} \\
+& \S{6}{110} + \S{7}{111}
\end{align*}
}{%
\visible<4->{%
\MedSkip{}
\begin{align*}
 \QState{}=& \alert<6>{\T{\alert<10>{2}}}  \\
 \TensOp{} & \alert<7>{\T{\alert<11>{1}}} \\
 \TensOp{} & \alert<8>{\T{\alert<12>{0}}}
\end{align*}}}
\end{frame}}


{
\def\TB#1#2{\LBP{}\ColorOne{\ensuremath{#1}}\ColorTwo{\ensuremath{#2}}}
\def\WO{\ColorOne{\ensuremath{w_1}}}\def\WT{\ColorTwo{\ensuremath{w_2}}}
\def\TE#1{\TwoSupOp{\QZero{}}{{\ExpPhase{2\pi \left(\LBP{}#1\right)}}\QOne{}}{+}}
\begin{frame}{Example using $n=2$ qubits}{Applying the theorem}
\Vskip{-4em}\TwoUnequalColumns{0.37\textwidth}{0.63\textwidth}{%
\ScrollProof{1}{4}{%
\Next{\Four}{\QState{} &= \SumPH{y}{2} \ExpPhase{2\pi wy}\ket{y} \\}
\Next{\Three}{w &= \LBP{}\WO{}\WT{} \\}
\gdef\CarryOne{\QState{} &= \SumPH{y}{2} \ExpPhase{2\pi \left(\LBP{}{\WO{}\WT{}}\right)y}\ket{y}}
\Last{\CarryOne}
}%
\only<5->{%
\begin{align*}
    \CarryOne
\end{align*}}
\only<5->{%

The exact possibilities for $w$:
\begin{center}
    \begin{tabular}{cc}
    \LBP{}\WO\WT & $w$ \\ \hline
    \visible<6->{\TB{0}{0} & \alert<10>{0} \\[0.2em]}%
    \visible<7->{\TB{0}{1} & \alert<11>{$\frac{1}{4}$}\\[0.2em]}%
    \visible<8->{\TB{1}{0} & \alert<10>{$\frac{1}{2}$} \\[0.2em]}%
    \visible<9->{\TB{1}{1} & \alert<11>{$\frac{3}{4}$}}
    \end{tabular}
\end{center}
}
}{%
\only<5->{%
\begin{align*}
= &\ColorThree{\TE{\WT}} \\ \TensOp{} &\ColorFour{\TE{\WO\WT}}
\end{align*}}
\only<10->{%
\begin{description}
    \item<10->[$\Hadamard(\ColorThree{top}) = \ket{0}$] then $\WT=0$, so \TB{w_1}{w_2}=\TB{0}{0}\mbox{ or }\TB{1}{0}. $\Hadamard(\ColorFour{bottom})$ thus yields $w=0$ or~$\frac{1}{2}$.
    \item<11->[$\Hadamard(\ColorThree{top}) = \ket{1}$] then $\WT=1$, so \TB{w_1}{w_2}=\TB{0}{1}\mbox{ or }\TB{1}{1}. We rotate away \WT{} and measure $\Hadamard(\ColorFour{bottom})$, yielding $w=\frac{1}{4}$ or~$\frac{3}{4}$.
\end{description}
But how do we \emph{rotate away}?
}
}
\end{frame}
}




\end{document}
